\section{Erd\H{o}s Problem \#776}

\subsection*{1) ROUND-2 OBJECTIVE}
\textbf{Path (C): obstruction/correction.} The round-1 writeup gave a valid but very coarse explicit upper bound. Since February 2026 there is a dedicated paper on this exact Erd\H{o}s--Trotter problem, so the right Round-2 objective is to replace the crude bound by the current strongest theorem and extract the asymptotic consequence. The problem is still not fully solved, but its first-order growth is now known.

\subsection*{2) Round-1 FOUNDATION USED}
I use the following Round-1 facts.

\begin{enumerate}
\item For $r>1$, an $r$-multiplicity antichain cannot realize $n-2$ distinct sizes.
\item Round-1 gave an explicit sufficient bound of the form
\[
 n_0(r)\le 8r+10,
\]
where $n_0(r)$ is the threshold beyond which $n-3$ distinct sizes are always achievable.
\end{enumerate}

\subsection*{3) NEW INSIGHT / TOOL (ROUND-2)}
The new input is the 2026 theorem of He and Tang. Write $g(n,r)$ for the maximum number of distinct set sizes appearing in an $r$-multiplicity antichain on $[n]$, and let $n_0(r)$ be the least integer such that
\[
 g(n,r)=n-3 \qquad \text{for every } n>n_0(r).
\]

\medskip
\noindent\textbf{Theorem 3.1 (He--Tang, 2026).}
\begin{enumerate}
\item $n_0(2)=3$ and $n_0(3)=8$.
\item For every integer $r\ge 4$,
\[
2r+2\le n_0(r)\le 2r+2\log_2 r+\log_2\log_2 r+15.
\]
In particular,
\[
 n_0(r)=2r+o(r).
\]
\end{enumerate}

\subsection*{4) ATTACK PLAN (ROUND-2)}
The exact remaining gap after Round-1 was quantitative: Round-1 only showed a coarse explicit upper bound, while the true threshold was unknown. The new plan is:

\begin{enumerate}
\item verify that the He--Tang theorem addresses the same threshold $n_0(r)$ as in the original problem;
\item import their current best upper and lower bounds;
\item derive the asymptotic corollary $n_0(r)=2r+o(r)$, which is a genuine first-order resolution of the threshold question.
\end{enumerate}

\subsection*{5) WORK (ROUND-2)}
There is one minor convention issue to settle first. The original Erd\H{o}s--Trotter formulation uses \emph{exactly} $r$ sets on each occurring level, whereas the Round-1 formulation used \emph{at least} $r$ sets. These conventions are equivalent for the extremal quantity $g(n,r)$: given an antichain with at least $r$ sets on each occurring level, one may keep an arbitrary subfamily of exactly $r$ sets from each occurring level. The antichain property is preserved and the set of occurring levels is unchanged.

Therefore Theorem 3.1 applies directly to the same threshold parameter in the original problem.

It follows immediately that the current quantitative status is:
\[
2r+2\le n_0(r)\le 2r+2\log_2 r+\log_2\log_2 r+15 \qquad (r\ge 4),
\]
and also
\[
 n_0(2)=3,\qquad n_0(3)=8.
\]

The most useful derived consequence is the first-order asymptotic.

\medskip
\noindent\textbf{Corollary 5.1.} \emph{As $r\to\infty$,}
\[
\frac{n_0(r)}{r}\to 2.
\]

\medskip
\noindent\emph{Proof.} Divide the bounds of Theorem 3.1 by $r$:
\[
2+\frac{2}{r}\le \frac{n_0(r)}{r}\le 2+\frac{2\log_2 r+\log_2\log_2 r+15}{r}.
\]
Since $(2\log_2 r+\log_2\log_2 r+15)/r\to 0$, the squeeze theorem gives $n_0(r)/r\to 2$. \hfill$\square$

Thus the Round-1 bound $n_0(r)\le 8r+10$ has now been superseded by a near-optimal estimate with the correct leading constant.

\subsection*{6) ADVERSARIAL VERIFICATION}
I check the main points where a mismatch could occur.

\begin{enumerate}
\item \emph{Does Theorem 3.1 use the same notion of multiplicity?} Yes, because the exact-$r$ and at-least-$r$ versions are equivalent for $g(n,r)$.

\item \emph{Is the asymptotic statement really justified?} Yes. The additive error in the upper bound is $2\log_2 r+O(\log_2\log_2 r)=o(r)$.

\item \emph{Does this fully solve Problem \#776?} No. The exact threshold for $r\ge 4$ is still unknown; the remaining gap is logarithmic.

\item \emph{Do the small cases fit the threshold definition with the strict inequality $n>n_0(r)$?} Yes. The theorem states the exact values of $n_0(2)$ and $n_0(3)$ under that convention.
\end{enumerate}

\subsection*{7) FINAL}
\textbf{UNRESOLVED (BUT STRICTLY ADVANCED)}

Problem \#776 is not fully closed, but the current state is far stronger than in Round-1: the exact small values are known,
\[
 n_0(2)=3,\qquad n_0(3)=8,
\]
and for general $r$ one has
\[
 n_0(r)=2r+o(r).
\]
Only a logarithmic gap remains between the best lower and upper bounds.

\subsection*{8) COMPLETION ESTIMATE (MANDATORY)}
\textbf{COMPLETION: 85\%}

\subsection*{9) REFERENCES}
[1] T. F. Bloom, \emph{Erd\H{o}s Problem \#776}, erdosproblems.com/776, accessed 2026-03-07.

[2] Y. He and Q. Tang, \emph{An Erd\H{o}s--Trotter problem on antichains with multiplicity $r$ on each occurring level}, arXiv:2602.09803, 2026.
